\chapter{统计}

\subsection*{总体和样本、数据获取、抽样方法}
\begin{enumerate}
    \item 某市为了分析全市 9800 名初中毕业生的数学考试成绩，共抽取 50 本试卷， 每本都是 30 份，则样本容量是\blkda{1500}.
    
    \item 从某年级 1000 名学生中抽取 125 名学生进行体重的统计分析, 就这个问题来说；“抽查的 125 名学生的体重” 是一个\blkda{样本}(选:总体，样本， 总体的容量，样本容量)

    \item 某中学高一年级有 400 人，高二年级有 320 人，高三年级有 280 人，若每人被抽到的可能性都为 0.2,用随机数表法在该中学抽取容量为 $n$ 的样本，则 $n$ 等于\blkda{$200$}.
    \begin{jieda}
        $0.2 \times (400+320+280) = 200$
    \end{jieda}

    \item 某学习小组在一次数学测验中，得 100 分的有 1 人，95 分的有 1 人，90 分的有 2 人，85 分的有 4 人，80 分和 75 分的各 1 人，则该小组成绩的平均分是\blkda{$87$}.

    \item 为了了解某市高三毕业生升学考试中数学成绩的情况，从参加考试的学生中随机抽查了1000 名学生的数学成绩进行统计分析,在这个问题中,下列说法正确的是\dotfill{(\kh{D})}
    \xx{总体指的是该市参加升学考试的全体学生}{个体指的是 1000 名学生中的每一名学生}{样本容量指的是 1000 名学生}{样本是指 1000 名学生的数学成绩}
    
    \item 已知一个总体含有 ${N}$ 个个体,要用简单随机抽样方法从中抽取一个个体,则在抽样过程中，每个个体被抽取的概率\dotfill{(\kh{C})}
    \xx{变小}{变大}{相等}{无法确定}
    \begin{jieda}
        一个总体含有 $N$ 个个体,要用简单随机抽样方法从中抽取一个个体，则在抽样过程中,每个个体被抽取的概率为 $\dfrac{1}{N}$ .
    \end{jieda}
    
    \item 现要完成下列2项抽样调查:
    \begin{dingautolist}{172}
        \item 从盒饼干中抽取 4 盒进行食品卫生检查
        \item 某中学共有 360 名教职工，其中一般教师 280 名，行政人员 55 名，后勤人员 25 名，为了了解教职工对学校在校务公开方面的意见，拟抽取一个容量为 72 的样本
    \end{dingautolist}
    较为合理的抽样方法是\dotfill{(\kh{A})}
    \xx{\ding{172}简单随机抽样， \ding{173} 分层抽样}{\ding{172}简单随机抽样， \ding{173} 简单随机抽样}{\ding{172}分层抽样， \ding{173} 分层抽样}{\ding{172}分层抽样， \ding{173} 简单随机抽样}
    \begin{jieda}
        \begin{dingautolist}{172}
            \item 总体中的个体数较少,宜用简单随机抽样；
            \item 总体是由差异明显的几部分组成,宜用分层抽样.
        \end{dingautolist}
    \end{jieda}
    \item 
    \begin{dingautolist}{172}
        \item 植物根据植株的高度及分枝部位等可以分为乔木、灌木和草木三大类，某植物园需要对其园中的不同植物的干重 (烘干后测定的质量) 进行测量
        \item 检测员拟对一批新生产的 1000 箱牛奶抽取 10 箱进行质量检测
    \end{dingautolist}
    上述两项调查应采用的抽样方法是\dotfill{(\kh{C})}
    \xx{\ding{172}简单随机抽样， \ding{173} 分层抽样}{\ding{172}简单随机抽样， \ding{173} 简单随机抽样}{\ding{172}分层抽样， \ding{173} 简单随机抽样}{\ding{172}分层抽样， \ding{173} 分层抽样}
    \begin{jieda}
        \begin{dingautolist}{172}
            \item 乔木、灌木、草木,分类明显,可以采用分层随机抽样
            \item 并未有明显分层特点,且样本容量较小,可以采用简单随机抽样；
        \end{dingautolist}
    \end{jieda}
    
    \item 一个总体分为 $A,B$ 两层，用分层抽样方法从总体中抽取一个容量为 20 的样本. 已知 $B$ 层中每个个体被抽到的概率都是 $\dfrac{1}{12}$ ,则总体中的个体数为\blkda{240}.
    \begin{jieda}
        因为用分层抽样方法从总体中抽取一个容量为 20 的样本. 由 $B$ 层中每个个体被抽到的概率都为 $\dfrac{1}{12}$ ,知道在抽样过程中每个个体被抽到的概率是 $\dfrac{1}{12}$，所以总体中的个体数为 ${20} \div  \dfrac{1}{12} = {240}$ .
    \end{jieda}

    \item 某高中的三个年级共有学生 1000 人，其中高一 300 人，高二 340 人，高三 360 人，该校现在要了解学生对校本课程的看法，准备从全校学生中抽取 50 人进行访谈，若采取分层抽样，且按年级来分层，则高一年级应抽取的人数是\blkda{15}.
    \begin{jieda}
        由题意可知抽样比为: $\dfrac{50}{1000} = \dfrac{1}{20}$，所以高一年级应抽取的人数是 ${300} \times  \dfrac{1}{20} = {15}$
    \end{jieda}

    \item 某校有学生 1485 人，教师 132 人，职工 33 人. 为有效防控甲型 H1N1 流感， 拟采用分层抽样的方法, 从以上人员中抽取 50 人进行相关检测, 则在学生中应抽取\blkda{45}人.
    \begin{jieda}
        $\dfrac{50}{1485+132+33} \times 1485 = 45$
    \end{jieda}

    \item 一个总体中的 80 个个体编号为 0,1,2,$\cdots$,79，并依次将其分为 8 个组，组号为 0,1,$\cdots$,7 , 要用下述抽样方法抽取一个容量为 8 的样本:即在第 0 组先随机抽取一个号码 $i$ ,则第 $k$ 组抽取的号码为 ${10k} + j$ ,其中 $j = \; \left\{  {\begin{array}{ll} i + k & i + k < {10} \\  i + k - {10} & i + k \geq  {10} \end{array}\;,\text{ 若先在 }0}\right.$ 组抽取的号码为6,则所抽到的8个号码依次为\blkda[5]{6, 17, 28, 39, 40, 51, 62, 73}.

    \item 从某地区 15000 位老人中随机抽取 500 人， 其生活能否自理的情况如下表所示:
    
    \begin{center}
    \begin{tabular}{|c|c|c|}
    \hline
    性别 & 男 & 女 \\
    \cline{1-3}
    能 & 178 & 278 \\
    \cline{1-3}
    不能 & 23 & 21 \\
    \cline{1-3}
    \hline
    \end{tabular}
    \end{center}
    
    则该地区生活不能自理的老人中男性比女性约多\blkda{60}人.
    \begin{jieda}
        $\dfrac{2}{500} \times 15000 = 60$
    \end{jieda}

    \item 一个总体分为 $A,B$ 两层,其个体数之比为 $4:1$,用分层抽样法从总体中抽取一个容量为 10 的样本，已知 $B$ 层中甲、乙都被抽到的概率为 $\dfrac{1}{28}$ ，则总体中的个体数是\blkda{40}.
    \begin{jieda}
        设$B$层中的个数为$n$，则$\dfrac{1}{C_n^2} = \dfrac{1}{28}$，$n = 8$，故总体有$8 \times (1+4) = 40$个.
    \end{jieda}

    \item 某公司甲、乙、丙、丁四个地区分别有 150 个、 120 个、 180 个、 150 个销售点. 公司为了调查产品销售的情况，需从这 600 个销售点中抽取一个容量为 100 的样本，记这项调查为 \ding{172} ；在丙地区中有 20 个特大型销售点，要从中抽取 7 个调查其收入和售后服务等情况，记这项调查为 \ding{173} . 则完成 \ding{172} 、 \ding{173} 这两项调查宜采用的抽样方法依次是\dotfill{(\kh{B})}
    \xx{分层抽样法, 系统抽样法}{分层抽样法, 简单随机抽样法}{系统抽样法, 分层抽样法}{简单随机抽样法, 分层抽样法}
    
    \item 将参加夏令营的 600 名学生编号为:001,002,$\cdots$, 600.采用系统抽样方法抽取一个容量为 50 的样本，且随机抽得的号码为 003 ，这 600 名学生分住在三个营区. 从 001 到 300 在第 I 营区，从 301 到 495 在第 II 营区，从 496 到 600 在第 III 营区. 三个营区被抽中的人数依次为\dotfill{(\kh{B})}
    \xx{26,16,8}{25,17,8}{25,16,9}{24,17,9}
    
    \item 某高级中学共有学生 2000 人，各年级男、 女生人数如表:
    \begin{center}
    \begin{tabular}{|c|c|c|c|}
    \hline
      & 高一 & 高二 & 高三 \\
    \cline{1-4}
    女生 & 373 & $x$ & $y$ \\
    \cline{1-4}
    男生 & 377 & 370 & $z$ \\
    \cline{1-4}
    \hline
    \end{tabular}
    \end{center}
    已知在全校学生中随机抽取 1 名，抽到高二年级女生的概率是 0.19 .
    \begin{enumerate}
        \item 现用分层抽样的方法在全校抽取 48 名学生，问应在高三年级抽取多少人?
        \item 已知 $y \geq  {245},z \geq  {245}$ ，求高三年级女生比男生多的概率.
    \end{enumerate}
    \begin{jieda}
        \begin{enumerate}
            \item 因为 $\dfrac{x}{2000} = {0.19}$ ,所以 $x = {380}$
            
            高三年级人数为 $y + z = {2000} - \left( {{373} + {377} + {380} + {370}}\right)  = {500}$，应在高三年级抽取 ${48} \times  \dfrac{500}{2000} = {12}$ 人
            
            \item 设高三年级女生比男生多的事件为A,高三年级女生、男生数记为$(y,z)$
            
            因为$y + z = {500}$ ,且 $y$ 、 $z$ 为正整数，基本事件有(245,255),(246,254),(247,253),$\cdots$,(255,245)共$11$个.
            
            事件A包含的基本事件有(251,249),(252,248),(253,247),(254,246),(255,245)共5个
            
            所以 $P\left( A\right)  = \dfrac{5}{11}$ .
        \end{enumerate}
    \end{jieda}
    
    \item 某校举行了 “环保知识竞赛”，为了了解本次竞赛成绩情况，从中随机抽取部分学生的成绩(得分均为整数，满分 100 分), 进行统计

    \begin{center}
    \begin{tabular}{|c|c|c|c|}
    \hline
    组号 & 分组 & 频数 & 频率 \\
    \cline{1-4}
    第 1 组 & [50,60) & 5 & 0.05 \\
    \cline{1-4}
    第 2 组 & [60,70) & $b$ & 0.35 \\
    \cline{1-4}
    第 3 组 & [70,80] & 30 & $c$ \\
    \cline{1-4}
    第 4 组 & [80,90] & 20 & 0.20 \\
    \cline{1-4}
    第 5 组 & [90,100) & 10 & 0.10 \\
    \cline{1-4}
    \multicolumn{2}{|c|}{合计} & $a$ & 1.00 \\
    \cline{1-4}
    \hline
    \end{tabular}
    
    \end{center}
    
    请根据频率分布表中所提供的数据，解答下列问题: 
    \begin{enumerate}
        \item 求 $a$ 、 $b$ 、 $c$ 的值及随机抽取一考生其成绩不低于 70 分的概率
        \item 若从成绩较好的 3、4、5 组中按分层抽样的方法抽取 6 人参加社区志愿者活动，并指定 2 名负责人，求从第 4 组抽取的学生中至少有一名是负责人的概率.
    \end{enumerate}
    \begin{jieda}
        \begin{enumerate}
            \item $c = 1 - \left( {{0.05} + {0.35} + {0.20} + {0.10}}\right)  =0.30$
            
            $a = \dfrac{5}{0.05} = 100$
            
            $b = {100} \times  {0.35} = {35},$
            
            由频率分布表可得成绩不低予 70(分)的概率为：$p = {0.30} + {0.20} + {0.10} = {0.60}$ .
            \item 若从成绩较好的3、4、5组中按分层抽样的方法抽取 6 人，则第3组抽取 $\dfrac{6}{60} \times  {30} = 3$ 人，第4组抽取 $\dfrac{6}{60} \times  {20} = 2$ 人，第5组抽取 $\dfrac{6}{60} \times  {10} = 1$ 人，从第 4 组抽取的学生中至少有一名是负责人的概率：$p = 1 - \dfrac{{C}_{4}^{2}}{{C}_{6}^{2}} = 1 - \dfrac{6}{15} = \dfrac{3}{5}.$
        \end{enumerate}
    \end{jieda}

    \item 从某校高三年级 800 名学生中随机抽取 50 名测量身高，据测量被抽取的学生的身高全部介于 155cm 和 195cm 之间，将测量结果按如下方式分成八组:第一组 $[155,160)$, 第二组$[160,165)$, $\cdots$, 第八组 $[190.195]$, 如图是按上述分组方法得到的频率分布直方图.
    
    \begin{center}
    \includegraphics[width=0.4\textwidth]{bodyxb/src/bo_d7ipq8491nqc7382lu50_3_216_1048_540_321_0.jpg}
    \end{center}
    
    \begin{enumerate}
        \item 根据已知条件填写下列表格:
        
        \begin{center}
        \begin{tabular}{|c|c|c|c|c|c|c|c|c|}
        \hline
        组别 & 一 & 二 & 三 & 四 & 五 & 六 & 七 & 八 \\
        \cline{1-9}
        样本数 & \blkda[1]{2} & \blkda[1]{4} & \blkda[1]{10} & \blkda[1]{10} & \blkda[1]{15} & \blkda[1]{4} & \blkda[1]{3} & \blkda[1]{2} \\
        \cline{1-9}
        \hline
        \end{tabular}
        \end{center}
        \item 试估计这所学校高三年级 800 名学生中身高在 180cm 以上(含 180cm)的人数为多少
        \item 在样本中，若第二组有 1 名男生，其余为女生，第七组有 1 名女生，其余为男生, 在第二组和第七组中各选一名同学组成实验小组, 问: 实验小组中恰有一男一女的概率是多少?
    \end{enumerate}
    \begin{jieda}
        \begin{enumerate}
            \item 如上
            \item $\dfrac{4+3+2}{50} \times 800 = 144$
            \item $P = \dfrac{1}{4} \times \dfrac{1}{3} + \dfrac{3}{4} \times \dfrac{2}{3} = \dfrac{7}{12}$
        \end{enumerate}
    \end{jieda}
\end{enumerate}

\clearpage
\subsection*{统计图表}

\begin{enumerate}
    \item 数据 21，19，31，25，28，18，30 的极差是\blkda{13}.
    \begin{jieda}
        数据 21,19,31,25,28,18,30 的极差是 ${31} - {18} = {13}$ .
    \end{jieda}
    
    \item 小明同学每天阅读数学文化相关的书籍，他每天阅读的页数分别为:4、5、4.5、5、 6、8、7、5、4.5、6(单位:页).下列图形中不利于描述这些数据的是\dotfill{(\kh{C})}
    \xx{条形图}{茎叶图}{散点图}{扇形图}
    \begin{jieda}
        条形图：是用宽度相同的条形的高度(或长度)表示数据的频数,故符合题意；
        
        茎叶图：即可以保留原始数据又可以方便记录数据,故符合题意；
        
        散点图: 用两组数据构成多个坐标点, 通常用于比较跨类别的成对数据, 不符合题意;
        
        扇形图: 是用整个圆表示总体, 用圆内各个扇形的大小表示各个部分占总体的百分数, 扇形图可以容易看出各个部分所占总体的比例, 故符合题意;
    \end{jieda}
    
    \item 下列关于散点图的说法中，正确的是 \dotfill{(\kh{B})}
    \xx{任意给定统计数据，都可以绘制散点图}{从散点图中可以看出两个量是否具有一定的关系}{从散点图中可以看出两个量的因果关系}{从散点图中无法看出数据的分布情况}
    \begin{jieda}
        散点图不适合用于展示百分比占比的数据, 另外数据量较少的数据也不适合用散点图表示,故 ${A}$ 错误；
        
        散点图能看出两个量是否具有一定关系,但是并一定是因果关系,故 ${B}$ 正确, ${C}$ 错误;
        
        散点图中能看出数据的分布情况,故 ${D}$ 错误.
    \end{jieda}

    \item 某同学将观察学校柚子树生长习性作为自主研究课题，他观察了校园内 6 株柚子树成熟结果个数 (两位数) 并用茎叶图 (如图所示) 做了记录，则这 6 株柚子树成熟结果个数的中位数为\blkda{$21.5$}.
    \begin{center}
    \includegraphics[width=0.15\textwidth]{bodyxb/src/bo_d7ipqc2lb0pc73f1h5bg_0_238_632_257_181_0.jpg}
    \end{center}

    
    \item 如图所示的是某班 60 名同学参加2011 年高中数学毕业会考所得成绩 (成绩均为整数)整理后画出的频率分布直方图，根据图中可得出的该班不及格(60 分以下)的同学的人数为\blkda{$15$}.
    \begin{center}
    \includegraphics[width=0.25\textwidth]{bodyxb/src/bo_d7ipqc2lb0pc73f1h5bg_0_1071_761_343_270_0.jpg}
    \end{center}
    \begin{jieda}
        由图可知: 分数是 60 以下的频率为 $\left( {{0.01} + {0.015}}\right)  \times  {10} = {0.25}$ ,所以人数为${60} \times  {0.25} = {15}$
    \end{jieda}

    \item 下列图形中具有相关关系的两个变量是\dotfill{(\kh{D})}
    \xx{\includegraphics[width=0.3\textwidth]{bodyxb/src/bo_d7ipqc2lb0pc73f1h5bg_0_243_1117_182_166_0.jpg}}{\includegraphics[width=0.3\textwidth]{bodyxb/src/bo_d7ipqc2lb0pc73f1h5bg_0_448_1119_178_163_0.jpg}}{\includegraphics[width=0.3\textwidth]{bodyxb/src/bo_d7ipqc2lb0pc73f1h5bg_0_657_1118_187_163_0.jpg}}{\includegraphics[width=0.3\textwidth]{bodyxb/src/bo_d7ipqc2lb0pc73f1h5bg_0_854_1119_177_163_0.jpg}}
    
    \item 下列四个图各反映两个变量的某种关系，其中可以看作具有较强线性相关关系的是\dotfill{(\kh{B})}
    \begin{center}
        \includegraphics[width=0.75\linewidth]{bodyxb//src/tongjitu.png}
    \end{center}

    \xx{\ding{172}  \ding{174}}{\ding{172}  \ding{175}}{\ding{173}  \ding{174}}{\ding{172}  \ding{173}  \ding{175}}
    
    \item 某大品牌家电公司从销售员工中随机抽出 50 名调查销售情况，销售额都在区间 $\left\lbrack  {5,{25}}\right\rbrack$ (单位:百万元)内,将其分成 5 组: $\lbrack 5,9),\lbrack 9,{13}),\lbrack {13},{17}),\lbrack {17},{21})$，$\left\lbrack  {{21},{25}}\right\rbrack$ ，并整理得到如右的频率分布直方图，据此估计销售员工销售额的平均值为\blkda{$14.52$}. (百万元) (同一组中的数据用该组区间的中点值为代表) .
    \begin{center}
    \includegraphics[width=0.4\textwidth]{bodyxb/src/bo_d7ipqc2lb0pc73f1h5bg_0_216_1661_486_309_0.jpg}
    \end{center}
    \begin{jieda}
        由题设 $\left( {{0.02} + a + {0.09} + {0.03} + {0.03}}\right)  \times  4 = 1$ ,可得 $a = {0.08}$，所以平均值为 ${0.08} \times  7 + {0.32} \times  {11} + {0.36} \times  {15} + {0.12} \times  {19} + {0.12} \times  {23} = {14.52}$ .
    \end{jieda}


    \item 已知甲，乙两组数据如茎叶图所示，其中 $m,n \in  \mathbf{N}$ ，若这两组数据的中位数相等，平均数也相等，则 $\dfrac{n}{m} =$ \blkda{$\dfrac{8}{3}$}.
    
    \begin{center}
    \includegraphics[width=0.25\linewidth]{bodyxb/src/bo_d7ipqc2lb0pc73f1h5bg_1_266_1232_289_211_0.jpg}
    \end{center}
    \begin{jieda}
        由茎叶图知甲的中位数为 ${30} + m$ ,乙的中位数为 $\dfrac{{32} + {34}}{2} = {33}$，由两组数据的中位数相同,则 ${30} + m = {33},m = 3$，由平均数相同,得 $\dfrac{{27} + {33} + {39}}{3} = \dfrac{{20} + n + {32} + {34} + {38}}{4}$ ,解得 $n = 8$，所以 $\dfrac{n}{m} = \dfrac{8}{3}$ .
    \end{jieda}


    \item 某学校组织全校学生参加网络安全知识竞赛，成绩(单位:分)的频率分布直方图如图所示，数据的分组依次为 $\lbrack {20},{40}),\lbrack {40},{60}),\lbrack {60},{80}),\left\lbrack  {{80},{100}}\right\rbrack$ ,若该校的学生总人数为 1000，则成绩不低于 60 分的学生人数为\blkda{$700$}.
    \begin{center}
    \includegraphics[width=0.4\textwidth]{bodyxb/src/bo_d7ipqc2lb0pc73f1h5bg_1_250_402_507_285_0.jpg}
    \end{center}
    \begin{jieda}
        根据频率分布直方图得该校的学生成绩不低于 60 分的学生人数为 ${1000} \times  \left( {{0.015} + {0.02}}\right)  \times  {20} = {700}$ .
    \end{jieda}

    \item 下列说法: \ding{172} 一组数据不可能有两个众数； \ding{173} 一组数据的方差必须是正数； \ding{174} 将一组数据中的每一个数据都加上或减去同一常数后，方差恒不变； \ding{175} 在频率分布直方图中， 每个小长方形的面积等于相应小组的频率.其中错误的有\blkda{\ding{172}\ding{173}}. (填写序号)
    

    
    \item 某校从参加语言测试的学生中随机抽取了 100 名，记录了他们的分数，将数据分成 6 组: $\lbrack {40},{50})$ ， $\lbrack {50},{60}),\lbrack {60},{70}),\lbrack {70},{80}),\lbrack {80},{90}),\left\lbrack  {{90},{100}}\right\rbrack$ ,并整理得到如下频率分布直方图. 若样本中分数低于 60 分的有 15 人，则图中数据 $a =$ \blkda{$0.01$}.
    
    \begin{center}
    \includegraphics[width=0.4\textwidth]{bodyxb/src/bo_d7ipqc2lb0pc73f1h5bg_1_216_857_475_365_0.jpg}
    \end{center}
    \begin{jieda}
        样本中分数低于 60 分的有 15 人,属于区间 $\lbrack {40},{50}),\lbrack {50},{60})$ ,由于学生中随机抽取了 100 名，因此分数在 $\lbrack {40},{50}),\lbrack {50},{60})$ 的频数为 15,因此这两个区间内的频率和为 $\dfrac{15}{100} = {0.15}$
        
        设区间 $\lbrack {40},{50})$ 的频率为 $x$ ,则 $\left\{  \begin{array}{l} x + {10a} = {0.15} \\  x + {10a} + {0.25} + {0.3} + {0.2} + {10a} = 1 \end{array}\right.$ ,解得 $a = {0.01}$ .
    \end{jieda}


    \item 已知甲、乙两组数据如茎叶图所示，其中 $m,n \in  \mathbf{N}$ . 若这两组数据的中位数相等，平均数也相等，则 $m \cdot  n =$ \blkda{$24$}.
    \begin{center}
    \includegraphics[width=0.2\textwidth]{bodyxb/src/bo_d7ipqc2lb0pc73f1h5bg_1_266_1232_289_211_0.jpg}
    \end{center}
    \begin{jieda}
        根据茎叶图知,乙的中位数是 $\dfrac{1}{2} \times  \left( {{32} + {34}}\right)  = {33}$ ,所以 $m = 3$，所以甲的平均数是 $\dfrac{1}{3} \times  \left( {{27} + {33} + {39}}\right)  = {33}$，所以 $\dfrac{1}{4} \times  \left( {{20} + n + {32} + {34} + {38}}\right)  = {33}$ ,解得 $n = 8$ .
    \end{jieda}

    \item 某学校为了了解学生的日平均睡眠时间(单位:$h$)，随机选择了 $n$ 名同学进行调查. 下表是这 $n$ 名同学的日睡眠时间的频率分布表.
    
    \begin{center}
    \begin{tabular}{|c|c|c|c|}
    \hline
    序号(i) & 分组(睡眠时间) & 频数(人数) & 频率 \\
    \cline{1-4}
    1 & [4,5) & 6 & 0.12 \\
    \cline{1-4}
    2 & [5,6] & \phantom{X} & 0.20 \\
    \cline{1-4}
    3 & [6,7) & $a$ & \phantom{X} \\
    \cline{1-4}
    4 & [7,8) & $b$ & \phantom{X} \\
    \cline{1-4}
    5 & [8,9) & \phantom{X} & 0.08 \\
    \cline{1-4}
    \hline
    \end{tabular}
    
    \end{center}
    \begin{enumerate}
        \item 求 $n$ 的值. 若 $a = {20}$ ，将表中数据补全
        \item 统计方法中，同一组数据常用该组区间的中点值(例如区间$[4,5)$的中点值是 $4.5$)作为代表. 若据此计算的上述数据的平均值为 $6.52$ , 求 $a$, $b$ 的值, 并由此估计该学校学生的日平均睡眠时间在 7 小时以上的概率.
    \end{enumerate}
    \begin{jieda}
        \begin{enumerate}
            \item  表格如下: 
            \begin{center}
            \begin{tabular}{|c|c|c|c|}
            \hline
            序号(i) & 分组(睡眠时间) & 频数(人数) & 频率 \\
            \cline{1-4}
            1 & [4,5) & 6 & 0.12 \\
            \cline{1-4}
            2 & [5,6] & 10 & 0.20 \\
            \cline{1-4}
            3 & [6,7) & 20 & 0.4 \\
            \cline{1-4}
            4 & [7,8) & 10 & 0.2 \\
            \cline{1-4}
            5 & [8,9) & 4 & 0.08 \\
            \cline{1-4}
            \hline
            \end{tabular}
            
            \end{center}
            \item 由题意可得 $a + b = {50} - 6 - {10} - 4 = {30}$，且 ${4.5} \times  {0.12} + {5.5} \times  {0.2} + {6.5} \times  \dfrac{a}{50} + {7.5} \times  \dfrac{b}{50} + {8.5} \times  {0.08} = {6.52}$，即 $a + b = {30}$ ,且 ${13a} + {15b} = {420}$ ,解得 $a = {15},b = {15}$.
            
            故学生的睡眠时间在7小时以上的频率等于 $\dfrac{15}{50} + {0.08} = {0.38}$ .
        \end{enumerate}
    \end{jieda}
    
    \item 空气质量指数 PM2.5 (单位: ${\mu {g}}/{{m}}^{3}$ ) 表示每立方米空气中可吸入肺颗粒物的含量， 这个值越高, 就代表空气污染越严重:
    
    \begin{center}
    \begin{tabular}{|c|c|c|c|c|c|c|}
    \hline
    PM2.5 日均浓度 & 0$\sim$ 35 & 35$\sim$ 75 & 75$\sim$ 115 & 115$\sim$ 150 & 150$\sim$ 250 & >250 \\
    \cline{1-7}
    空气质量级别 & 一级 & 二级 & 三级 & 四级 & 五级 & 六级 \\
    \cline{1-7}
    空气质量类别 & 优 & 良 & 轻度污染 & 中度污染 & 重度污染 & 严重污染 \\
    \cline{1-7}
    \hline
    \end{tabular}
    \end{center}

    甲、乙两城市 2013 年 2 月份中的 15 天对空气质量指数 PM2.5 进行监测, 获得 PM2.5 日均浓度指数数据如茎叶图所示:
    \begin{center}
        \includegraphics[width=0.25\linewidth]{bodyxb//src/tongjitu2.png}
    \end{center}

    \begin{enumerate}
        \item 根据你所学的统计知识估计甲、乙两城市 15 天内哪个城市空气质量总体较好？并说明理由.
        \item 在 15 天内任取 1 天，估计甲、乙两城市空气质量类别均为优或良的概率
    \end{enumerate}
    \begin{jieda}
        \begin{enumerate}
            \item 由茎叶图知：甲城市的中位数为：61,乙城市的中位数为 79，并且甲的大多集中在 65 以下, 乙的大多集中在 76 以上，所以甲城市空气质量总体较好;
            
            \item 甲城市空气质量类别为优或良的有 10 天，所以甲城市空气质量类别优或良的概率为 $P = \dfrac{10}{15} = \dfrac{2}{3}$，乙城市空气质量类别为优或良的有 5 天，所以乙城市空气质量类别优或良的概率为 $P = \dfrac{5}{15} = \dfrac{1}{3}$，所以甲、乙两城市空气质量类别均为优或良的概率为 $P = \dfrac{2}{3} \times  \dfrac{1}{3} = \dfrac{2}{9}$ .
        \end{enumerate}
    \end{jieda}
    
    \item 某校高二年级半期考试后，为了解本次考试的情况，在整个年级中随机抽取了 200 名学生的数学成绩，将成绩分为 $\lbrack {30},{50}),\lbrack {50},{70}),\lbrack {70},{90}),\lbrack {90},{110}),\lbrack {110},{130}),\left\lbrack  {{130},{150}}\right\rbrack$ , 共 6 组, 得到如图所示的频率分布直方图.
    \begin{center}
    \includegraphics[width=0.4\textwidth]{bodyxb/src/bo_d7ipqc2lb0pc73f1h5bg_2_229_1323_530_411_0.jpg}
    \end{center}

    \begin{enumerate}
        \item 求实数 $a$ 的值
        \item 在样本中, 采取按比例分层抽样的方法从成绩在 $[90,150]$ 内的学生中抽取 13 名, 问其中成绩在 $[130,150]$ 的学生有几名?
        \item 根据图中的样本数据，假设同组中每个数据用该组区间的中点值代替，试估计本次考试的平均分.
    \end{enumerate}
    \begin{jieda}
        \begin{enumerate}
            \item 由频率分布直方图知：
            $\left( {{0.0025} + {0.005} \times  2 + {0.01} + a + {0.015}}\right)  \times  {20} = 1$ ,
            解得 $a = {0.0125}$ .
            \item 采取分层抽样,[130,150]的学生个数为： $\dfrac{0.005}{{0.005} + {0.015} + {0.0125}} \times  {13} = 2$ ,
            即成绩在 $\left\lbrack  {{130},{150}}\right\rbrack$ 的学生有 2 名.
            \item 由频率分布直方图知：平均数为：
            $\left( {{40} \times  {0.0025} + {60} \times  {0.005} + {80} \times  {0.01} + {100} \times  {0.015} + {120} \times  {0.0125} + {140} \times  {0.005}}\right)  \times  {20} = {98}$
        \end{enumerate}
    \end{jieda}
    
    \item 某校为了提高学生的反诈骗意识，举办了反诈骗知识竞赛，从所有答卷中随机抽取 100 份作为样本，将样本的成绩(满分 100 分，成绩均为不低于 40 分的整数)分成六组: $\lbrack {40},{50})$ ， $\lbrack {50},{60}),\lbrack {60},{70}),\lbrack {70},{80}),\lbrack {80},{90}),\left\lbrack  {{90},{100}}\right\rbrack  ,$ 得到如图所示的频率分布直方图.

    \begin{center}
    \includegraphics[width=0.4\textwidth]{bodyxb/src/bo_d7ipqc2lb0pc73f1h5bg_3_241_203_526_381_0.jpg}
    \end{center}
    \begin{enumerate}
        \item 求频率分布直方图中 $a$ 的值与样本成绩的平均数
        \item 用分层随机抽样的方法从 $\left\lbrack  {{70},{80}}\right) ,\left\lbrack  {{90},{100}}\right\rbrack$ 两个区间共抽取出 8 名学生，再从这 8 名学生中随机抽取 2 名依次进行交流分享,求第二个交流分享的学生成绩在区间 $\lbrack {70},{80})$ 的概率
        \item 学校决定从知识竞赛中抽出成绩最好的两个同学甲乙进行现场知识抢答赛，比赛共设三个项目，每个项目胜方得 1 分，负方得 0 分，没有平局，三个项目比赛结束后，总得分高的人获得冠军. 已知甲在三个项目中获胜的概率分别为 $\dfrac{1}{2}$ ， $\dfrac{2}{5}$ ， $p$ ，各项目的比赛结果相互独立,甲至少得 1 分的概率是 $\dfrac{41}{50}$ ,甲乙两人谁获得最终胜利的可能性大? 并说明理由.
    \end{enumerate}
    \begin{jieda}
        \begin{enumerate}
            \item (1)由频率之和为 1 ,结合频率分布直方图可得${10} \times  \left( {{0.005} + {0.01} + {0.02} + a + {0.025} + {0.01}}\right)  = 1$ ,解得 $a = {0.03}$
            
            样本成绩的平均数约为 $\bar{x} = {45} \times  {0.05} + {55} \times  {0.1} + {65} \times  {0.2} + {75} \times  {0.3} + {85} \times  {0.25} + {95} \times  {0.1} = {74}$ .
            
            \item 由频率分布直方图知,样本答卷成绩在 $\lbrack {70},{80}),\left\lbrack  {{90},{100}}\right\rbrack$ 的学生比为 $3 : 1$ ,用分层抽样的方法从在 $\lbrack {70},{80})$ 的学生中抽取 $\dfrac{3}{4} \times  8 = 6$ (人)，从 $\left\lbrack  {{90},{100}}\right\rbrack$ 的学生中抽取 $\dfrac{1}{4} \times  8 = 2$ (人).
            从这 8 名学生中随机抽取 2 名依次进行交流分享,有 $8 \times  7 = {56}$ 种情形,
            第一个和第二个交流分享的学生成绩均在区间 $\lbrack {70},{80})$ ,有 $6 \times  5 = {30}$ 种情形,
            第一个交流分享的学生成绩在区间 $\left\lbrack  {{90},{100}}\right\rbrack$ ,第二个交流分享的学生成绩在区间 $\lbrack {70},{80})$ ,有
            $2 \times  6 = {12}$ 种情形
            
            故所求概率为 $P = \dfrac{{30} + {12}}{56} = \dfrac{3}{4}$ .
            \item 乙最终获胜的可能性大. 理由如下：
            由题意,甲至少得 1 分的概率是 $\dfrac{41}{50}$ ,
            可得 $1 - \left( {1 - \dfrac{1}{2}}\right) \left( {1 - \dfrac{2}{5}}\right) \left( {1 - p}\right)  = \dfrac{41}{50}$ ,其中 $0 \leq  p \leq  1$ ,解得 $p = \dfrac{2}{5}$ ,
            则甲得 2 分或 3 分的概率为:
            $P = \dfrac{1}{2} \times  \dfrac{2}{5} \times  \left( {1 - \dfrac{2}{5}}\right)  + \dfrac{1}{2} \times  \left( {1 - \dfrac{2}{5}}\right)  \times  \dfrac{2}{5} + \left( {1 - \dfrac{1}{2}}\right)  \times  \dfrac{2}{5} \times  \dfrac{2}{5} + \dfrac{1}{2} \times  \dfrac{2}{5} \times  \dfrac{2}{5} = \dfrac{2}{5}.$
            所以乙得分为 2 分或 3 分的概率为 $\dfrac{3}{5}$ ,乙最终获胜的可能性大.
        \end{enumerate}
    \end{jieda}
\end{enumerate}
\clearpage
\subsection*{统计估计}

\begin{enumerate}
    \item 从一堆苹果中任取 5 只, 称得它们的质量分别为 125、123、121、124、127 (单位:克)，则估计每个苹果的平均质量为\blkda{124}(克)，该堆苹果的标准差 $S =$ \blkda{2}(克).(用数字作答)
    \begin{jieda}
        $\bar{x} = \dfrac{1}{5}\left( {{125} + {124} + {121} + {123} + {127}}\right) = {124}$
        
        ${S}^{2} = \dfrac{1}{5}\lbrack {\left( {125} - {124}\right) }^{2} + {\left( {124} - {124}\right) }^{2} + {\left( {121} - {124}\right) }^{2} + {\left( {123} - {124}\right) }^{2} + {\left( {127} - {124}\right) }^{2}\rbrack = 4$
        
        $S = \sqrt{4} = 2$ .
    \end{jieda}
    
    \begin{center}
    \includegraphics[width=0.3\textwidth]{bodyxb/src/bo_d7ipq9qlb0pc73f1h570_0_225_526_423_289_0.jpg}
    \end{center}


    \item 为了调查某班学生做数学题的基本能力, 随机抽查了部分学生的某次数学考试成绩(这次考试试卷满分为 100 分), 这些分数的分组区间为[45,55), [55,65), [65,75), [75,85), [85,95], 这些分数的分组区间为[45,55), [55,65), [65,75), [75,85), [85,95], 由此得到频率分布直方图如图所示, 则这些学生的平均分为\blkda{64}.
    
    \item 已知某企业职工年收入的频率分布如表所示试估计该企业职工的平均年收入为\blkda{5.1}. (万元).
    \begin{center}
    \begin{tabular}{|c|c|}
    \hline
    年收入范围 (万元) & 频率 $f$ \\
    \cline{1-2}
    [3,5) & 0.60 \\
    \cline{1-2}
    [5,7) & 0.25 \\
    \cline{1-2}
    [7,9) & 0. 15 \\
    \cline{1-2}
    \hline
    \end{tabular}
    \end{center}
    
    \item 从某公司抽取 25 名雇员，考察他们的工资情况如下表所示:
    
    \begin{center}
    \begin{tabular}{|c|c|c|c|c|c|c|c|c|}
    \hline
    年薪(千元) & 135.0 & 95.0 & 80.0 & 70.0 & 60.0 & 52.0 & 40.0 & 31.0 \\
    \cline{1-9}
    人数 & 1 & 1 & 2 & 1 & 3 & 4 & 1 & 12 \\
    \cline{1-9}
    \hline
    \end{tabular}
    \end{center}
    
    估计该公司职工年薪的中位数是\blkda{40.0}(千元). (结果精确到 0.1)
    
    估计该公司职工年薪的标准差是\blkda{25.5}(千元).(结果精确到 0.1)

    \item 某学校为了解该校 600 名男生的百米成绩(单位:$s$)，随机选择了 50 名学生进行调查, 下图是这 50 名学生百米成绩的频率分布直方图. 根据样本的频率分布，估计这 600 名学生中成绩在 $\left\lbrack  {{13},{15}}\right\rbrack$ (单位: ${s}$ ) 内的人数大约是\blkda{120}.
    
    \begin{center}
    \includegraphics[width=0.2\textwidth]{bodyxb/src/bo_d7ipq9qlb0pc73f1h570_0_229_1667_305_418_0.jpg}
    \end{center}
    \begin{jieda}
        $(0.02 \times 1 + 0.18 \times 1) \times 600 = 120$
    \end{jieda}

    \item 给出下列 4 个命题，其中正确命题的序号是\blkda{\ding{172}\ding{173}\ding{174}}.
    \begin{dingautolist}{172}
        
        \item 在大量的试验中，事件 $A$ 出现的频率可以作为事件 $A$ 出现的概率的估计值
        \item 样本标准差 $S = \sqrt{\dfrac{{\left( {x}_{1} - \bar{x}\right) }^{2} + {\left( {x}_{2} - \bar{x}\right) }^{2} + \cdots  + {\left( {x}_{n} - \bar{x}\right) }^{2}}{n - 1}}\left( {n \geq  2}\right)$ 可以作为总体标准差的点估计值
        \item 随机抽样就是使得总体中每一个个体都有同样的可能性被选入样本的一种抽样方法
        \item 分层抽样就是把总体分成若干部分，然后在每个部分指定某些个体作为样本的一种抽样方法

    \end{dingautolist}
    
    \item 小波一星期的总开支分布如图(1)所示，一星期的食品开支如图(2)所示，则小波一星期的鸡蛋开支占总开支的百分比为\dotfill{(\kh{C})}
    \xx{$30\%$}{$10\%$}{$3\%$}{不能确定}

    \twotu{\includegraphics[width=0.6\textwidth]{bodyxb/src/bo_d7ipq9qlb0pc73f1h570_1_222_936_495_308_0.jpg}}{\includegraphics[width=0.6\textwidth]{bodyxb/src/bo_d7ipq9qlb0pc73f1h570_1_759_941_463_276_0.jpg}}
    
    \item 某地 2004 年第一季度应聘和招聘人数排行榜前 5 个行业的情况列表如下
    
    \begin{center}
        \includegraphics[width=0.5\linewidth]{bodyxb//src/448cd12932678ea7c0be0886fe3e8eac.jpg}
    \end{center}
    
    若用同一行业中应聘人数与招聘人数比值的大小来衡量该行业的就业情况，则根据表中数据，就业形势一定是\dotfill{(\kh{B})}
    \xx{计算机行业好于化工行业}{建筑行业好于物流行业}{机械行业最紧张}{营销行业比贸易行业紧张}
    
    \item 在发生某公共卫生事件期间，有专业机构认为该事件在一段时间没有发生在规模群体感染的标志为 “连续 10 天，每天新增疑似病例不超过 7 人”。根据过去 10 天甲、乙、 丙、丁四地新增疑似病例数据，一定符合该标志的是\dotfill{(\kh{D})}
    \xx{甲地:总体均值为 3 ，中位数为 4}{乙地:总体均值为 1 ，总体方差大于 0}{丙地: 中位数为 2 , 众数为 3}{丁地:总体均值为 2，总体方差为 3}
    
    \item 为了普及环保知识，增强环保意识，某大学随机抽取 30 名学生参加环保知识测试，测试成绩(单位:分)如图所示，假设得分值的中位数为 ${m}_{e}$ ,众数为 ${m}_{o}$ ,平均值为 $\bar{x}$ ,则\dotfill{(\kh{D})}
    
    \begin{center}
    \includegraphics[width=0.4\textwidth]{bodyxb/src/bo_d7ipq9qlb0pc73f1h570_2_211_522_493_286_0.jpg}
    \end{center}

    \xx{${m}_{e} = {m}_{o} = \bar{x}$}{${m}_{e} = {m}_{o} < \bar{x}$}{${m}_{e} < {m}_{o} < \bar{x}$}{${m}_{o} < {m}_{e} < \bar{x}$}
    
    \item 某市电视台为宣传在伦敦举办的 “2012 年奥运会”，随机对该市 ${15} \sim  {65}$ 岁的人群抽取了 $n$ 人，回答问题 “2012 年奥运会是第几届奥运会？”，统计结果如图表所示，右侧为不同年龄段抽样的频率分布直方图:
    
    \begin{minipage}{0.45\textwidth}
        \centering
        \begin{tabular}{|c|c|p{2cm}|p{2cm}|}
        \hline
        组号 & 分组 & 回答正确的人数 & 正确人数占本组的概率 \\
        \cline{1-4}
        第 1 组 & [15,25) & 5 & 0.5 \\
        \cline{1-4}
        第 2 组 & [25,35) & $a$ & 0.9 \\
        \cline{1-4}
        第 3 组 & [35,45) & 27 & $x$ \\
        \cline{1-4}
        第 4 组 & [45,55) & $b$ & 0.36 \\
        \cline{1-4}
        第 5 组 & [55,65) & 3 & $y$ \\
        \cline{1-4}
        \hline
        \end{tabular}
    \end{minipage}
    \hfill
    \begin{minipage}{0.45\textwidth}
        \centering
        \includegraphics[width=0.9\textwidth]{bodyxb/src/bo_d7ipq9qlb0pc73f1h570_2_936_1089_472_422_0.jpg}
    \end{minipage}
    
    \begin{enumerate}
        \item 分别求出 $a,b,x,y$ 的值
        \item 在第 2，3，4 组回答正确的人中用分层抽样的方法抽取 6 人，则第 2，3，4 组每组应各抽取多少人?
        \item 在(2)的前提下，电视台决定在所抽取的 6 人中随机抽取 2 人颁发幸运奖，求所抽取的人中第 2 组至少有 1 人获得幸运奖的概率.
    \end{enumerate}
    \begin{jieda}
        \begin{enumerate}
            \item $a = 18,b = 9,x = 0.9, y = 0.2$
            \item 第2,3,4组回答正确的共有54人.
            $\therefore$ 利用分层抽样在54人中抽取6人,
            每组分别抽取的人数为: 第2组: $\dfrac{18}{54} \times  6 = 2$ 人,第3组: $\dfrac{27}{54} \times  6 = 3$ 人,
            第4组: $\dfrac{9}{54} \times  6 = 1$ 人.
            \item 设第 2 组的 2 人为 ${A}_{1},{A}_{2}$ ,第 3 组的 3 人为 ${B}_{1},{B}_{2}$，${B}_{3}$ ,第 4 组的 1 人为 ${C}_{1}$ ,则从 6 人中抽取 2 人的所有可能的结果为 $\left( {{A}_{1},{A}_{2}}\right) ,\left( {{A}_{1},{B}_{1}}\right) ,\left( {{A}_{1},{B}_{2}}\right) ,\left( {{A}_{1},{B}_{3}}\right) ,\left( {A}_{1}, {C}_{1}\right) ,\left( {{A}_{2},{B}_{1}}\right) ,\left( {{A}_{2},{B}_{2}}\right) ,\left( {{A}_{2},{B}_{3}}\right) ,\left( {{A}_{2},{C}_{1}}\right) ,\left( {{B}_{1},{B}_{2}}\right)$ ,
            
            $\left( {{B}_{1},{B}_{3}}\right) ,\left( {{B}_{1},{C}_{1}}\right) ,\left( {{B}_{2},{B}_{3}}\right) ,\left( {{B}_{2},{C}_{1}}\right) ,\left( {{B}_{3},{C}_{1}}\right)$ ,共 15种情况,其中每 2 组至少有 1 人被抽中的情况有 $\left( {A}_{1}\right.$ ,$\left. {A}_{2}\right) ,\left( {{A}_{1},{B}_{1}}\right) ,\left( {{A}_{1},{B}_{2}}\right) ,\left( {{A}_{1},{B}_{3}}\right) ,\left( {{A}_{1},{C}_{1}}\right) ,\left( {{A}_{2},{B}_{1}}\right)$ ,$\left( {{A}_{2},{B}_{2}}\right) ,\left( {{A}_{2},{B}_{3}}\right) ,\left( {{A}_{2},{C}_{1}}\right)$ ,共 9 种情况.
            
            故第 2 组至少有 1 人获得幸运奖的概率为 $\dfrac{9}{15} = \dfrac{3}{5}$ .
        \end{enumerate}
    \end{jieda}
    
    \item 某中学共有 1000 名学生参加了该地区高三第一次质量检测的数学考试，数学成绩如下表所示:
    
    \begin{center}
    \begin{tabular}{|c|c|c|c|c|c|}
    \hline
    数学成绩分组 & $\lbrack 0,{30})$ & $\lbrack {30},{60})$ & [60, 90) & [90,120) & [120,150] \\
    \cline{1-6}
    人数 & 60 & 90 & 300 & $x$ & 160 \\
    \cline{1-6}
    \hline
    \end{tabular}
    \end{center}
    \begin{enumerate}
        \item 为了了解同学们前段时间的复习情况，学校将采用分层抽样的方法抽取 100 名同学进行问卷调查，甲同学在本次测试中数学成绩为 95 分，求他被抽中的概率
        \item 已知本次数学成绩的优秀线为 110 分，试根据所提供数据估计该中学达到优秀线的人数
        \item 作出频率分布直方图，并估计该学校本次考试的数学平均分.
        
        \begin{center}
        \includegraphics[width=0.3\textwidth]{bodyxb/src/bo_d7ipq9qlb0pc73f1h570_3_195_728_564_494_0.jpg}
        \end{center}
    \end{enumerate}
    \begin{jieda}
        \begin{enumerate}
            \item $p = \dfrac{1}{10}$
            \item 290
            \item 估计该学校本次考试的数学平均分为90分
        \end{enumerate}
    \end{jieda}
    
    \item 某班级共有 60 名学生，先用抽签法从中抽取部分学生调查他们的学习情况， 若每位学生被抽到的概率为 $\dfrac{1}{6}$ . 
    \begin{enumerate}
        \item 求从中抽取的学生数
        \item 若抽查结果如下表
        
        \begin{center}
        \begin{tabular}{|c|c|c|c|c|}
        \hline
        每周学习时间(小时) & [0,10) & [10,20) & [20,30) & [30,40] \\
        \cline{1-5}
        人数 & 2 & 4 & $x$ & 1 \\
        \cline{1-5}
        \hline
        \end{tabular}
        \end{center}
        
        先确定 $x$ ，再完成频率分布直方图
        \begin{center}
        \includegraphics[width=0.4\textwidth]{bodyxb/src/bo_d7ipq9qlb0pc73f1h570_3_238_1647_534_408_0.jpg}
        \end{center}
        \item 估计该班学生每周学习时间的平均数
    \end{enumerate}
    \begin{jieda}
        \begin{enumerate}
            \item 10人
            \item $x = 3$
            \item 18小时
        \end{enumerate}
    \end{jieda}

\end{enumerate}

\subsection*{补充:抽样方法}
\begin{enumerate}
    \item  某学校共有教师 490 人，其中不到 40 岁的有 350 人，40 岁及以上的有 140 人，为了检查普通话在该校教师中的推广普及情况，用分层抽样的方法，从全体教师中抽取一个容量为 70 的样本进行普通话水平测试, 其中在不到 40 岁的教师中应抽取的人数是\blkda{50}.

    \item 网络上流行一种 “ ${QQ}$ 农场游戏”，这种游戏通过虚拟软件模拟种植与收获的过程. 为了了解本班学生对此游戏的态度，高三(6)班计划在全班 60 人中展开调查, 根据调查结果, 班主任计划采用系统抽样的方法抽取若干名学生进行座谈, 为此先对 60 名学生进行编号为:01,02,03,$\cdots$60，已知抽取的学生中最小的两个编号为 03,09 ，则抽取的学生中最大的编号为\blkda{57}.
    
    \item 某班级有 50 名学生，现要采取系统抽样的方法在这 50 名学生中抽出 10 名学生,将这 50 名学生随机编号 1 $\sim$  50 号,并分组,第一组 1 $\sim$  5 号,第二组 6 $\sim$  10 号, $\cdots \cdots$ ,第十组 ${46} \sim  {50}$ 号,若在第三组中抽得号码为 12 的学生,则在第八组中抽得号码为\blkda{37}的学生.
    
    \begin{center}
    \includegraphics[width=0.3\textwidth]{bodyxb/src/bo_d7ipqhqlb0pc73f1h5dg_0_233_1159_460_189_0.jpg}
    \end{center}
    
    \item 某单位 200 名职工的年龄分布情况如图, 现要从中抽取 40 名职工作样本. 用系统抽样法, 将全体职工随机按 1$\sim$ 200 编号，并按编号顺序平均分为 40 组(1 $\sim$  5 号，6 $\sim$  10 号，$\cdots$，196 $\sim$  200 号). 若第 5 组抽出的号码为 22，则第 8 组抽出的号码应是\blkda{37}. 若用分层抽样方法，则 40 岁以下年龄段应抽取\blkda{20}人.

    \item 某企业共有职工 150 人，其中高级职称 15 人，中级职称 45 人，初级职称 90 人. 现采用分层抽样抽取容量为 30 的样本，则抽取的各职称的人数分别为\blkda{$3, 9, 18$}.

    \item 某高中在校学生有 2000 人. 为了响应 “阳光体育运动” 的号召, 学校开展了跑步和登山比赛活动. 每人都参与而且只参与其中一项比赛，各年级参与比赛的人数情况如下表:
    
    \begin{center}
    \begin{tabular}{|c|c|c|c|}
    \hline
    \phantom{X} & 高一年级 & 高二年级 & 高三年级 \\
    \cline{1-4}
    跑步 & a & $b$ & $C$ \\
    \cline{1-4}
    登山 & $x$ & $y$ & $z$ \\
    \cline{1-4}
    \hline
    \end{tabular}
    \end{center}
    
    其中 $a : b : c = 2 : 3 : 5$ ，全校参与登山的人数占总人数的 $\dfrac{2}{5}$ . 为了了解学生对本次活动的满意程度，从中抽取一个 200 人的样本进行调查，则从高二年级参与跑步的学生中应抽取\blkda{36}.

    \item 某中学采用系统抽样方法，从该校高一年级全体 800 名学生中抽 50 名学生做牙齿健康检查. 现将 800 名学生从 1 到 800 进行编号. 已知从 33 $\sim$  48 这 16 个数中取的数是 39，则在第 1 小组 1 $\sim$  16 中随机抽到的数是\blkda{7}.

    \item \begin{enumerate}
        \item 某学校为了了解 2012 年高考数学的考试成绩, 在高考后对 1200 名学生进行抽样调查, 其中文科 400 名考生, 理科 600 名考生, 艺术和体育类考生共 200 名，从中抽取 120 名考生作为样本.
        \item 从 10 名家长中抽取 3 名参加座谈会. 
    \end{enumerate}
    I . 简单随机抽样法 II .系统抽样法 III.分层抽样法.
    
    问题与方法配对正确的是\dotfill{(\kh{A})}
    \xx{(1)III, (2)I}{(1)I , (2)II}{(1) II , (2) III}{(1) III, (2) II}
    
    \item 将参加夏令营的 600 名学生编号为:001,002$\cdots$, 600. 采用系统抽样方法抽取一个容量为 50 的样本, 且随机抽得的号码为 003 , 这 600 名学生分住在三个营区. 从 001 到 300 在第 I 营区, 从 301 到 495 在第 II 营区, 从 496 到 600 在第 III 营区. 三个营区被抽中的人数依次为\dotfill{(\kh{B})}
    \xx{26,16,8}{25,17,8}{25,16,9}{24,17,9}
    
    \item 在 100 个零件中，有一级品 20 个，二级品 30 个，三级品 50 个，从中抽取 20 个作为样本.
    \begin{dingautolist}{172}
        \item 采用随机抽样法:抽签取出 20 个样本
        \item 采用系统抽样法:将零件编号为 00,01,$\cdots$,99 , 然后平均分组抽取 20 个样本
        \item 采用分层抽样法:从一级品，二级品，三级品中抽取 20 个样本
    \end{dingautolist}
    下列说法中正确的是\dotfill{(\kh{A})}
    \xx{无论采用哪种方法，这 100 个零件中每一个被抽到的概率都相等}{\ding{172}\ding{173}两种抽样方法，这 100 个零件中每一个被抽到的概率都相等\ding{174}并非如此}{\ding{172}\ding{174}两种抽样方法，这 100 个零件中每一个被抽到的概率都相等； \ding{173} 并非如此}{采用不同的抽样方法，这 100 个零件中每一个零件被抽到的概率是各不相同的}
    
    \item 一个城市有 210 家百货商店, 其中大型商店有 20 家, 中型商店有 40 家, 小型商店有 150 家. 为了掌握各商店的营业情况, 要从中抽取一个容量为 21 的样本，按分层抽样方法抽取样本时，各类百货商店要分别抽取多少家？写出抽样过程.
    \begin{jieda}
        \begin{dingautolist}{172}
            \item 样本容量与总体中的个体数的比值为 $\dfrac{21}{210} = \dfrac{1}{10}$
            
            \item 确定要抽取的各种商店的数目,大型商店为$20 \times \dfrac{1}{10}$ = 2 (家),中型商店为$40 \times \dfrac{1}{10}$ = 4 (家),小型商店为$150 \times \dfrac{1}{10}$ = 15 (家)；
            
            \item 采用简单随机抽样的方法在各层中分别抽取大型商店2家、中型商店4家、小型商店15家,这样便得到了所要抽取的样本.
        \end{dingautolist}
    \end{jieda}
    
    \item 某单位有 2000 名职工，老年、中年、青年分布在管理、技术开发、营销、生产各部门中，如下表所示:
    
    \begin{center}
    \begin{tabular}{|c|c|c|c|c|c|}
    \hline
    人数 & 管理 & 技术开发 & 营销 & 生产 & 共计 \\
    \cline{1-6}
    老年 & 40 & 40 & 40 & 80 & 200 \\
    \cline{1-6}
    中年 & 80 & 120 & 160 & 240 & 600 \\
    \cline{1-6}
    青年 & 40 & 160 & 280 & 720 & 1200 \\
    \cline{1-6}
    共计 & 160 & 320 & 480 & 1040 & 2000 \\
    \cline{1-6}
    \hline
    \end{tabular}
    \end{center}
    \begin{enumerate}
        \item 若要抽取 40 人调查身体状况，则应怎样抽样？
        \item 若要开一个 25 人的讨论单位发展与薪金调整方面的座谈会，则应怎样抽选出席人?
        \item 若要抽 20 人调查对某运动会筹备情况的了解，则应怎样抽样？
    \end{enumerate}
    \begin{jieda}
        \begin{enumerate}
            \item 用分层抽样,并按老年4人,中年12人,青年24人抽取.
            \item 用分层抽样,并按管理2人,技术开发4人,营销6人,生产13人抽取.
            \item 用系统抽样. 对全部2 000人随机编号,号码从$0001 \sim 2000$,每100号分为一组,从第一组中用随机抽样抽取一个号码,然后将这个号码分别加100,200,$\cdots$,1900,共20人组成一个样本.
        \end{enumerate}
    \end{jieda}
    
    \item 某公路设计院有工程师 6 人，技术员 12 人，技工 18 人，要从这些人中抽取 $n$ 个人参加市里召开的科学技术大会. 如果采用系统抽样和分层抽样的方法抽取, 不用剔除个体，如果参会人数增加 1 个，则在采用系统抽样时，需要在总体中先剔除 1 个个体,求 $n$ .
    \begin{jieda}
        由题意知采用系统抽样和分层抽样方法抽取,不用剔除个体；如果样本容量增加一个,则在采用系统抽样时,需要在总体中先剔除1个个体
                
        $\because$ 总体容量为 $6 + {12} + {18} = {36}$ .
        $\therefore$ 当样本容量是 $n$ 时,由题意知,系统抽样的间隔为 $\dfrac{36}{n}$ ,分层抽样比例为 $\dfrac{n}{36}$ ,抽取的工程师人数为 $\dfrac{n}{36} \times  6 = \dfrac{n}{6}$ ,技工人数为 $\dfrac{n}{36} \times  {18} = \dfrac{n}{2}$ ,
        $\because n$ 应是6的倍数,36的约数,
        $\therefore n = 6,{12},{18}$ ,
        若当样本容量为 $n + 1$ 时,总体容量是35人,系统抽样的间隔为 $\dfrac{35}{n + 1}$ ,
        $\because \dfrac{35}{n + 1}$ 必须是整数,
        $\therefore n$ 只能取 6 ,即 $n = 6$ .
    \end{jieda}
    
    \item 某地区有小学 21 所, 中学 14 所, 大学 7 所, 现采取分层抽样的方法从这些学校中抽取 6 所学校对学生进行视力调查
    \begin{enumerate}
        \item 求应从小学、中学、大学中分别抽取的学校数目
        \item 若从抽取的 6 所学校中随机抽取 2 所学校做进一步数据分析
        \begin{dingautolist}{172}
            \item 列出所有可能的抽取结果
            \item 求抽取的 2 所学校均为小学的概率
        \end{dingautolist}
    \end{enumerate}
    \begin{jieda}
        \begin{enumerate}
            \item 从小学、中学、大学中分别抽取的学校数目为3、2、1.
            \item \begin{dingautolist}{172}
                \item 在抽取到的6所学校中,3所小学分别记为A1,A2,A3,2所中学分别记为A4,A5,大学记为A6,则抽取2所学校的所有可能结果为
                $\{ {A}_{1},{A}_{2}\} ,\{ {A}_{1},{A}_{3}\} ,\{ {A}_{1},{A}_{4}\} ,\{ {A}_{1},{A}_{5}\} ,$
                
                $\{ {A}_{1},{A}_{6}\} ,\{ {A}_{2},{A}_{3}\} ,\{ {A}_{2},{A}_{4}\} ,\{ {A}_{2},{A}_{5}\} ,\{ {A}_{2},{A}_{6}\} ,\{ {A}_{3},{A}_{4}\} ,\{ {A}_{3},{A}_{5}\} ,\{ {A}_{3},{A}_{6}\} ,$
                
                $\{ {A}_{4},{A}_{5}\} ,\{ {A}_{4},{A}_{6}\} ,\{ {A}_{5},{A}_{6}\}.$ 共15种
                \item 从6所学校中抽取的2所学校均为小学(记为事件B)的所有可能结果为\{A1,A2\},\{A1,A3\},\{A2,A3\},共3种，所以 $P\left( B\right)  = \dfrac{3}{15} = \dfrac{1}{5}$ .
            \end{dingautolist}
        \end{enumerate}
    \end{jieda}
\end{enumerate}

\clearpage

\subsection*{专题: 百分位数}
\begin{enumerate}
    \item 现有一组数据:3,4,6,8,9,10,12,13，其第 70 百分位数为\blkda{10}.
    \begin{jieda}
        因为 $8 \times  {70}\%  = {5.6}$,从小到大排列数据:3,4,6,8,9,10,12,13,向上取整取第 6 位数为 10 .
    \end{jieda}
    
    \item 现有 7 张卡片，分别写上数字 2, 4, 5, 5, 6, 9, 16，则这 7 个数的第 75 百分位数是\blkda{9}.
    \begin{jieda}
        由于 $7 \times  {75}\%  = {5.25}$,则这 7 个数的第 75 百分位数是 9
    \end{jieda}
    \item 某工厂选取 40 名工人，将他们随机分成两组，每组 20 人，第一组工人用第一种生产方式，第二组工人用第二种生产方式. 根据工人完成生产任务的工作时间(单位:分钟)绘制了如下茎叶图:
    \twotu{\includegraphics[width=0.7\textwidth]{bodyxb/src/bo_d7j0uhilb0pc73f1l7sg_0_255_1348_559_222_0.jpg}}{\includegraphics[width=0.7\textwidth]{bodyxb/src/bo_d7j0uhilb0pc73f1l7sg_0_863_1200_531_378_0.jpg}}
    \begin{enumerate}
        \item 求 40 名工人完成生产任务所需时间的第 75 百分数
        \item 为了解该工厂职工的基本信息，从工厂中抽取了 100 个职工的体重数据，发现全部介于 45 公斤到 75 公斤之间，现将 100 个体重数据分为 6 组:第一组 $\lbrack {45},{50})$，第二组 $\lbrack {50},{55})$， ...，第六组 $\lbrack {70},{75})$，得到如图 2 所示的频率分布直方图. 其中第一组有 2 人，第二组有 13 人. 求 $\dfrac{b}{c}$与 $a$的值.
    \end{enumerate}
    \begin{jieda}
        \begin{enumerate}
            \item 40 名工人完成生产任务所需时间按从小到大排列为:
            
            61,61,62,63,63,65,65,67,68,69,70,70,71,72,72,72,72,74,75,77,
            
            78,81,82,82,83,84,84,84,87,87,90,90,91,91,91,92,92,93,94,94,因为 ${40} \times  {75}\%  = {30}$,
            
            所以第 75 百分数为 $\dfrac{{87} + {90}}{2} = {88.5}$
            \item 依题意 $\left( {{2a} + {2b} + {2c}}\right)  \times  5 = 1$，则 $a + b + c = {0.1}$，
            
            又因为 ${5a} \times  {100} = 2$，所以 $a = {0.004}$，
            
            因为 ${5b} \times  {100} = {13}$,所以 $b = {0.026}$,
            
            所以 $c = {0.07}$,
            
            所以 $\dfrac{b}{c} = \dfrac{13}{35},a = {0.004}$.
        \end{enumerate}
    \end{jieda}
    \item 为加强学生睡眠监测督导, 学校对高中三个年级学生的日均睡眠时间进行调查.根据分层随机抽样法, 学校在高一、高二和高三年级中共抽取了 100 名学生的日均睡眠时间作为样本，其中高一 35 人，高二 33 人. 已知该校高三年级一共 512 人.
    \begin{enumerate}
        \item 学校高中三个年级一共有多少个学生？
        \item 若抽取 100 名学生的样本极差为 2，数据如下表所示(其中 $x < {10}$， $n$是正整数)
        
        \begin{center}
        \begin{tabular}{|c|c|c|c|c|c|}
        \hline
        日均睡眠时间(小时) & $x$& 8.5 & 9 & 9.5 & 10 \\
        \cline{1-6}
        学生数量 & $n$& 32 & 13 & 11 & 4 \\
        \cline{1-6}
        \hline
        \end{tabular}
        \end{center}
        
        求该样本的第 40 百分位数.
    \end{enumerate}
    
    \begin{jieda}
        \begin{enumerate}
            \item 设学校高中三个年级一共有 $N$个学生,
            
            因为采用分层抽样法抽取一个容量为 100 的样本,
            
            在高一年级抽取了 35 人，高二年级抽取了 33 人，
            
            所以高三抽取的人数为:100-35-33=32人，
            
            又因为高三年级一共 512 人，所以有: $\dfrac{32}{512} = \dfrac{100}{N}$，解得 $N = {1600}$.
            
            所以学校高中三个年级一共有 1600 个学生.
            \item 因为抽取 100 名学生的样本极差为 $2,x < {10}$，
            
            所以 ${10} - x = 2 \Rightarrow  x = 8$， $n = {100} - {32} - {13} - {11} - 4 = {40}$.
            
            又因为 ${100} \times  {0.4} = {40}$,
            
            所以样本的第 40 百分位数为: $\dfrac{8 + {8.5}}{2} = {8.25}$(小时).
        \end{enumerate}
    \end{jieda}

    \item 某学校为了获得该校全体高中学生的体育锻炼情况，按男、女学生的比例分别抽样调查了 48 名男生和 27 名女生的每周锻炼时间. 通过计算得到男生每周锻炼时间的平均数为 7.6 小时, 方差为 7; 女生每周锻炼时间的平均数为 6.4 小时, 方差为 8 .
    \begin{enumerate}
        \item 若该校男生总数为 1280 ，求该校学生总数
        \item 若所选 27 名女生每周锻炼时间从小到大排列后的第 9 至第 13 个数据依次为 5、5.3、5.6、5.8、5.9，求所选女生样本的第 40 百分位数
        \item 求所有样本数据的平均数和方差(精确到 0.001)
    \end{enumerate}
    \begin{jieda}
        \begin{enumerate}
            \item 设该校学生总数为 $n$,依题意, $\dfrac{1280}{n} = \dfrac{48}{{48} + {27}}$,解得 $n = {2000}$,所以该校学生总数为 2000 .
            \item 由 ${27} \times  {40}\%  = {10.8}$，得所选女生样本的第 40 百分位数为第 11 个数 5.6.
            \item 所有样本数据的平均数 $\bar{x} = \dfrac{{48} \times  {7.6} + {27} \times  {6.4}}{{48} + {27}} = {7.168}$;
            
            所有样本数据的方差为 ${s}^{2} = \dfrac{48}{{48} + {27}}\left\lbrack  {7 + {\left( {7.6} - {7.168}\right) }^{2}}\right\rbrack   + \dfrac{27}{{48} + {27}}\left\lbrack  {8 + {\left( {6.4} - {7.168}\right) }^{2}}\right\rbrack   \approx  {7.692}$.
        \end{enumerate}
    \end{jieda}
\end{enumerate}